% !TeX root = ta.tex
% ==========================================
% BAB VII PENUTUP
% ==========================================
\cleardoublepage
\chapter{PENUTUP}
\label{chap:penutup}
\section{Kesimpulan}
\begin{enumerate}
	\item Tugas akhir ini membandingkan tiga arsitektur (CLM, \textit{Explainable} CNN, EfficientNetB0) dan menemukan performa bergantung pada kualitas representasi fitur. CLM (ekstraksi fitur \textit{non-trainable}) paling rentan dan mengalami kolaps kelas. \textit{Explainable} CNN (dilatih dari nol) mengalami penurunan \textit{specificity} drastis pada beberapa skenario. EfficientNetB0 (dengan \textit{transfer learning}) mencatat akurasi \textit{single-split} tertinggi (96,63\%), namun turun menjadi 83,5\% $\pm$ 5,11\% saat divalidasi menggunakan \textit{Grouped Stratified 5-Fold Cross-Validation}.
	
	\item Integrasi Grad-CAM pada EfficientNetB0 (model performa terbaik) belum mampu menyoroti area tumor secara relevan, baik kuantitatif maupun kualitatif. Hal ini disebabkan tiga faktor yang saling memperkuat, yaitu keterbatasan resolusi spasial Grad-CAM akibat \textit{pooling} berulang, ketajaman citra yang berbeda antar kelas, dan ketergantungan Grad-CAM pada arsitektur \textit{lightweight} yang kurang presisi dalam lokalisasi. Temuan ini mengindikasikan adanya \textit{trade-off}, yaitu akurasi tertinggi justru menghasilkan penjelasan visual yang paling kurang relevan, sehingga objektivitas model belum dapat divalidasi secara klinis hanya berdasarkan Grad-CAM, dan diperlukan uji kausal serta metode XAI lain untuk memastikannya.
\end{enumerate}

\section{Saran}
Berdasarkan keterbatasan yang ditemukan selama tugas akhir ini, beberapa saran diajukan untuk pengembangan lebih lanjut.

\begin{enumerate}
	\item Memperbanyak jumlah data pada \textit{dataset} yang digunakan, misalnya dengan menjalin kerja sama untuk memperoleh data citra dari rumah sakit tertentu. Dengan demikian, \textit{dataset} yang diperoleh dapat lebih spesifik dan sesuai dengan karakteristik data di rumah sakit tersebut. Model yang dihasilkan pun dapat lebih mendekati kebutuhan diagnosis klinis yang sebenarnya.
	
	\item Mengintegrasikan model EfficientNetB0 yang telah dibuat pada tugas akhir ini dengan model segmentasi. Dengan demikian, sistem tidak hanya mampu menyatakan ada atau tidaknya tumor, tetapi juga dapat menunjukkan batas dan area tumor secara lebih presisi. Integrasi ini sekaligus dapat menjadi pembanding bagi hasil visualisasi Grad-CAM yang masih kurang akurat pada tugas akhir ini.
	
	\item Menyesuaikan proses pengolahan data maupun pengaturan \textit{hyperparameter} dengan karakteristik masing-masing model, mengingat setiap arsitektur memiliki kebutuhan yang berbeda. Pada EfficientNetB0 misalnya, seberapa besar bagian model yang ikut dilatih ulang (\textit{fine-tuning}) turut memengaruhi hasil, namun hal ini belum ikut diuji secara sistematis pada tugas akhir ini sehingga perlu menjadi perhatian pada penelitian selanjutnya.
	
	\item Melibatkan pakar biomedis seperti dokter spesialis radiologi untuk mengevaluasi secara langsung hasil prediksi maupun visualisasi Grad-CAM yang dihasilkan model. Hal ini penting mengingat penilaian yang dilakukan pada tugas akhir ini masih bersifat kuantitatif dan belum divalidasi dari sudut pandang klinis yang sesungguhnya.
	
	\item Mengeksplorasi kemungkinan pengubahan citra \textit{CT Scan} menjadi citra MRI sintesis. Perlu dicatat bahwa pendekatan ini belum ditujukan khusus untuk pendeteksian tumor sehingga efektivitasnya masih harus dikaji lebih lanjut. Dengan cara ini, keterbatasan detail jaringan lunak pada \textit{CT Scan} berpotensi dapat diatasi tanpa bergantung pada ketersediaan mesin MRI yang masih terbatas. Pendekatan ini juga berpotensi memperjelas lokalisasi tumor pada hasil \textit{explainability} model.
\end{enumerate}